\documentclass[11pt]{article}
\usepackage[margin=1in]{geometry}
\usepackage{lmodern,microtype}
\title{Q21 Comprehension Check}
\date{}
\begin{document}
\maketitle
\textbf{Question:} A colleague proposes a warehouse robot built as a 4-stage pipeline: object detector $\to$ grasp planner $\to$ motion controller $\to$ safety monitor. Each stage is a separately trained network. They report that each stage is ``95\% accurate in unit tests'' and therefore expect roughly 95\% end-to-end success. (a) What is wrong with the arithmetic, and under what assumption would 95\% actually be the right expectation? (b) Name two composition costs their design incurs that a monolithic end-to-end policy would not. (c) Their next proposal is to make the stages talk by passing natural-language descriptions between them, ``so we can read the logs.'' What does this interface choice cost them relative to structured embeddings?

\vspace{1em}
\textbf{Answer:} (a) Simple multiplication of per-stage accuracies ($0.95^4 \approx 0.81$) is only correct if stage errors are independent; in practice errors correlate (hard inputs fail at multiple stages), so end-to-end success is usually \emph{lower} than $0.81$ --- compounding error, not the unit-test product. (b) Any two of: compounding/propagated error across fixed seams with no way for downstream stages to correct upstream mistakes; interface brittleness (retraining the detector silently shifts the grasp planner's input distribution); cross-component credit assignment (a failed grasp cannot be localized to detector vs.\ planner without extra instrumentation); added serial latency. (c) Natural language is the lowest-bandwidth, lossiest interface: continuous geometry (poses, grasp points) cannot round-trip through text without quantization error, gradients cannot flow across the seam for joint fine-tuning, and generation adds latency --- they gain log readability (observability) but pay in accuracy ceiling and trainability of the joints.
\end{document}
